% Shared body for both layouts. Floats use the \widefig / \widetab
% environments, which expand to figure*/table* in two-column mode and to
% figure/table in one-column mode, so a single file serves both builds.
% ════════════════════════════════════════════════════════════════════════
\section{Introduction}
\label{sec:intro}

Retrieval-augmented generation has become the default way to put a language model
in front of a document corpus \citep{lewis2020rag,gao2023ragsurvey}. The standard
recipe embeds passages, retrieves the nearest neighbours of an embedded question,
and lets the model answer from what came back. It works well when the answer is
contained in a passage that resembles the question.

A large class of enterprise questions is not like that. Consider a question a
financial analyst would actually ask: \emph{for every company where a particular
person is an officer or director, list that company's disclosed subsidiaries and
summarize the cybersecurity risks it reported.} Three facts are required. The
person's role is disclosed in an insider-ownership filing. The subsidiaries are
disclosed in an exhibit attached to the annual report. The risks are narrative
prose inside a numbered section of that same report. These are three separate
documents, in three formats, produced by three different processes.

The difficulty is not that the passages are hard to rank. It is that no passage
contains the combination. We measured this directly on our corpus: the surname of
the subject of our hardest question appears in \textbf{0 of 5{,}210} extracted
risk-factor sections, and only \textbf{1 of 151} subsidiary names appears anywhere
in its own parent's risk-factor section. Ownership and risk are, empirically,
disjoint. \Cref{fig:containers} shows the structure that produces this.

This has a consequence that is easy to state and easy to overlook: a better
encoder cannot fix it, and neither can a longer context window. Retrieval can
only return text that exists. If the conjunction of facts a question needs is
never co-located in any passage, then ranking passages more accurately retrieves
the same insufficient evidence. The ceiling is a property of the corpus, not of
the retriever.

There is a second and subtler failure. Similarity search has no mechanism for
filtering by entity, because ``belongs to company $X$'' is not a semantic property
of a passage. Two risk-factor sections about supply-chain disruption are near each
other in embedding space whether or not they belong to the company being asked
about. At small scale this rarely bites, because few candidates compete. We show
that it bites hard at 508{,}714 chunks drawn from 5{,}210 filers, and --- more
usefully --- that it bites on exactly the questions where similarity search is
supposed to be the right tool.

The obvious alternative is to store the structure explicitly: resolve entities to
stable identifiers, materialize the relationships as edges, and attach a
provenance key to every edge so that any claim can be traced to the document that
supports it. This is the family of designs now loosely called GraphRAG
\citep{edge2024graphrag,peng2024graphragsurvey}. The question this paper addresses
is not whether that is possible but whether it is \emph{worth it}, and on which
question shapes --- measured against two baselines that were given every
reasonable advantage.

We are deliberately narrow. This is an empirical systems study on one corpus with
one model and twenty questions. It contributes measurements and diagnosed failure
modes, not a new algorithm. Our contributions are:

\begin{itemize}
  \item \textbf{A corpus and benchmark with hard entity identity.} Twelve monthly
        EDGAR feeds ingested into a relational store and a graph
        (\Cref{tab:corpus}), with 20 questions in four categories. The multi-hop
        subjects were selected on measured evidence: the surname must appear in
        \emph{zero} risk-factor sections and must be shared by two or more
        distinct filer identifiers, so that resolving identity by name is
        demonstrably wrong rather than merely inelegant
        (\Cref{sec:design}).
  \item \textbf{A matched three-way comparison.} 180 programmatically scored
        attempts under one model at temperature 0, with the fan-out guardrail and
        repair policy applied identically to both query-writing arms
        (\Cref{sec:arch}), and with two of the four categories chosen so that the
        graph should lose.
  \item \textbf{Two results that contradict our expectation.} Vector RAG lost the
        category it was expected to win, and plain text-to-SQL beat GraphRAG on
        cross-document entity resolution. We report both with the same prominence
        as the wins and diagnose each (\Cref{sec:failures}).
  \item \textbf{A cost accounting.} GraphRAG's accuracy is bought:
        780{,}500 tokens against text-to-SQL's 158{,}171, and a repair call
        triggered on 34 of 60 attempts against 7 of 60
        (\Cref{tab:main}, \Cref{tab:repair}). We also report the reverse cost ---
        text-to-SQL's generated joins averaged 92\,s and peaked at 445\,s on the
        multi-hop category.
  \item \textbf{Eight defects found in the measurement harness itself.} Each
        would have produced a confident and wrong conclusion had it gone
        unnoticed, including one ground-truth error found after the results were
        first computed (\Cref{sec:repro}). We argue that an unaudited benchmark
        harness is itself a leading source of false findings.
\end{itemize}

\Cref{sec:related} places the work; \Cref{sec:dataset} and \Cref{sec:arch}
describe the corpus and the three architectures; \Cref{sec:design} gives the
protocol; \Cref{sec:results} reports what happened; \Cref{sec:failures} explains
the failures; \Cref{sec:threats} states what these measurements cannot support.


\begin{widefig}[htbp]
\centering
\begin{tikzpicture}[
  font=\footnotesize,
  >={Stealth[length=1.7mm,width=1.3mm]},
  doc/.style={draw=gray!55, rounded corners=1.5pt, fill=white, align=center,
              font=\scriptsize, inner sep=3pt, text width=3.9cm, minimum height=1.0cm},
  emb/.style={doc, draw=veccol!70, fill=veccol!10, line width=0.7pt},
  noemb/.style={doc, draw=warncol!75, fill=warncol!10, line width=0.7pt},
  sep/.style={doc, draw=sqlcol!70, fill=sqlcol!10, line width=0.7pt},
  wrap/.style={draw=gray!45, dashed, rounded corners=3pt, inner sep=6pt},
  edge/.style={draw=graphcol!75, rounded corners=1.5pt, fill=graphcol!12,
               align=center, font=\tiny, inner sep=2.5pt, text width=3.3cm,
               minimum height=0.66cm, line width=0.6pt},
  ann/.style={align=center, font=\tiny, text width=5.2cm, inner sep=1pt},
  arr/.style={->, draw=gray!60, line width=0.5pt},
]
%% ── column 1: the containers, each annotated directly BELOW itself so that
%%    nothing can collide with the boxes or the arrows leaving them ─────────
\node[emb]   (i1a) at (0,0) {\textbf{Item 1A} Risk Factors\\[1pt]
                             \tiny narrative prose, mean 116{,}779 chars};
\node[ann, below=1.5mm of i1a] (a1)
  {\textcolor{veccol}{\textbf{embedded}} $\rightarrow$ 508{,}714 chunks\\
   \textcolor{gray!70}{the whole universe similarity search can reach}};

\node[noemb, below=3.0mm of a1] (ex21) {\textbf{EX-21} Subsidiaries\\[1pt]
                             \tiny structured list: name + jurisdiction};
\node[ann, below=1.5mm of ex21] (a2)
  {\textcolor{warncol}{\textbf{not embedded}} $\rightarrow$ 212{,}124 rows\\
   \textcolor{gray!70}{a list does not benefit from an encoder}};

\node[wrap, fit=(i1a)(a1)(ex21)(a2),
      label={[font=\scriptsize\bfseries, text=gray!80, inner sep=2pt]above:%
      10-K annual report \textmd{\tiny(one accession number)}}] (tenk) {};

\node[sep, below=7mm of tenk.south] (f345) {\textbf{Form 3 / 4 / 5}\\[1pt]
                             \tiny a different filing entirely};
\node[ann, below=1.5mm of f345] (a3)
  {\textcolor{sqlcol}{\textbf{not embedded}} $\rightarrow$ 150{,}917 rows\\
   \textcolor{gray!70}{carries the CIK that separates people sharing a surname}};

%% ── column 2: the graph projection ───────────────────────────────────────
\node[edge] (e1) at (7.6,0.0)   {\texttt{:FILED}\\\tiny 167{,}050};
\node[edge] (e2) at (7.6,-1.5)  {\texttt{:SUBSIDIARY\_OF}\\\tiny 212{,}124};
\node[edge] (e3) at (7.6,-3.0)  {\texttt{:OFFICER\_OF} / \texttt{:DIRECTOR\_OF}\\\tiny 86{,}180 / 70{,}167};
\node[wrap, draw=graphcol!55, fit=(e1)(e2)(e3),
      label={[font=\scriptsize\bfseries, text=graphcol, inner sep=2pt]above:%
      graph projection \textmd{\tiny(accession\_number on every edge)}}] (g) {};

\draw[arr] (i1a.east)  -- ++(0.5,0) |- (e1.west);
\draw[arr] (ex21.east) -- ++(0.5,0) |- (e2.west);
\draw[arr] (f345.east) -- ++(0.5,0) |- (e3.west);

%% ── the measured disjointness ────────────────────────────────────────────
\node[draw=badcol!70, fill=badcol!8, line width=0.7pt, rounded corners=2pt,
      align=left, font=\scriptsize, inner sep=5pt, text width=12.4cm,
      anchor=north west] at ($(tenk.north west)+(0,-7.55)$)
  {\textbf{Why no single chunk answers a multi-hop question.} Measured on this
   corpus: the subject's surname appears in \textbf{0 of 5{,}210} Item~1A
   sections, and only \textbf{1 of 151} subsidiary names appears anywhere in its
   own parent's Item~1A. Role, ownership and risk are written in three containers
   by three teams, so the combination exists in no retrievable passage.};
\end{tikzpicture}

\caption{The three facts a multi-hop question needs live in three different
containers. Item~1A is narrative prose and is the only section embedded; EX-21 is
a structured list attached to the same filing; insider roles are in a separate
filing type entirely. Every projected edge carries the accession number of the
filing that asserts it.}
\label{fig:containers}
\end{widefig}

% ════════════════════════════════════════════════════════════════════════
\section{Related Work}
\label{sec:related}

\textbf{Retrieval-augmented generation.} Grounding a generator in retrieved
passages was formalized by \citet{lewis2020rag} and rests on dense retrieval
\citep{karpukhin2020dpr} and sentence-level encoders \citep{reimers2019sbert}.
The literature since is large \citep{gao2023ragsurvey}, and much of it optimizes
the ranking step: better encoders, rerankers, hybrid sparse-dense scoring, query
rewriting. Our negative result on similarity search is orthogonal to that line of
work. We do not claim the retriever ranked badly; we claim that on questions whose
answer spans documents, no ranking of passages can return evidence that no passage
contains, and that entity membership is not a property the embedding encodes.

\textbf{Structured querying as a baseline.} Translating questions into executable
queries has its own long line, from \citet{zhong2017seq2sql} to the cross-domain
Spider benchmark \citep{yu2018spider}. We treat text-to-SQL not as a competitor to
be defeated but as the strongest available baseline, and we give it the complete
schema including the tables that make the multi-hop answers reachable by joins.
This matters for the interpretation of our results: text-to-SQL wins one of the
four categories outright and beats our graph implementation in a second.

\textbf{Knowledge graphs and graph-based retrieval.} Knowledge graphs
\citep{hogan2021kg} predate the current interest, and ``GraphRAG'' is now an
overloaded label \citep{peng2024graphragsurvey}. The best-known formulation
\citep{edge2024graphrag} builds a community structure over an LLM-extracted graph
and is aimed at query-focused summarization --- global questions such as ``what
are the themes in this corpus''. Our design targets a different shape. The graph
is built from filing metadata and structured exhibits rather than by LLM
extraction; the traversal resolves \emph{which entities} are in scope, and the
answer text is then fetched by primary key from the relational store rather than
summarized from the graph. Only the first hop is LLM-authored; expansion is
deterministic, for reasons we quantify in \Cref{sec:failures}.

\textbf{Multi-hop question answering.} HotpotQA \citep{ho2020hotpotqa} and MuSiQue
\citep{trivedi2022musique} established multi-hop evaluation, generally over
Wikipedia, where the hops are between topically related articles and entity names
are largely unambiguous. Our corpus differs in the property that drives our main
result: entity \emph{identity} is genuinely ambiguous. Three distinct people in
this corpus share the surname of our featured subject, and only a stable
identifier separates them. A benchmark where names resolve cleanly cannot expose
the failure mode we observe.

\textbf{Evaluation.} Automated RAG evaluation increasingly uses an LLM judge
\citep{es2024ragas}. We deliberately do not, because the object under study is
whether a system states falsehoods, and hallucination in the judge
\citep{ji2023hallucination} would be indistinguishable from hallucination in the
system. Our scoring is a deterministic function of required entities, forbidden
entities, keyword presence and citation correctness (\Cref{sec:design}). This
buys reproducibility at the cost of rewarding only the answers we anticipated.

Finally, we position this paper honestly: it is a single-corpus empirical study
whose value is in its measurements and its diagnosed failures, not in a new
retrieval algorithm.


% ════════════════════════════════════════════════════════════════════════
\section{Corpus Construction}
\label{sec:dataset}

\subsection{Source and document structure}

The corpus is drawn from the U.S. Securities and Exchange Commission's EDGAR
system \citep{sec2024edgar}, which publishes every mandatory disclosure filed by
every U.S. public company. We ingested twelve monthly XBRL feeds labelled
2025-09 through 2026-08. Filing dates in the resulting corpus span 2025-08 to
2026-08, because a monthly feed includes filings dated at the end of the preceding
month; \Cref{fig:months} shows the distribution, with the small 2025-08 bucket
being that spillover.

The structural point that motivates this paper is that an annual report is not one
document. A Form 10-K is a container: a mandated sequence of numbered
\emph{Items}, plus attached \emph{Exhibits}, filed under a single accession
number. Three containers matter here.

\textbf{Item 1A, Risk Factors}, is narrative prose in which a company enumerates
what could go wrong. It is the longest narrative section, averaging 116{,}779
characters and reaching 1{,}060{,}461 in this corpus (\Cref{tab:items}). It is the
only section we embed.

\textbf{Exhibit 21} is an attachment, not a section, required by Item 601(b)(21)
of Regulation S-K. It lists subsidiaries and their jurisdiction of incorporation,
so it parses into rows rather than prose. Two properties matter. First, the
distribution is extremely skewed: 275 filers disclose a single subsidiary while
62 disclose more than 500, and the largest discloses 2{,}780
(\Cref{fig:subdist}). Second, and more consequentially,
\emph{EX-21 is incomplete by design}: rule 601(b)(21)(ii) permits a filer to omit
any subsidiary that is not a ``significant subsidiary''. The exhibit is therefore
a floor on corporate structure and never a census, and the graph inherits that
limit. An absent subsidiary is not evidence that none exists.

\textbf{Forms 3, 4 and 5} are separate filings, not part of the 10-K. Officers,
directors and ten-percent owners report their positions and transactions in
structured XML. This is the container that makes identity resolvable, because each
filing carries the insider's own Central Index Key (CIK) --- a stable SEC
identifier. Two directors who share a surname are separable on CIK even though
their names are not.

\Cref{fig:containers} summarizes the arrangement and the measured disjointness
that follows from it.

\subsection{Ingestion}

The pipeline runs in nine stages: fetch and parse the monthly feeds; fetch the
documents each filing manifests; extract narrative Items from HTML; parse EX-21
and the ownership forms; stage to object storage; load the relational store;
embed; and build indexes and the graph projection. \Cref{fig:pipeline} gives the
wall-clock breakdown, totalling roughly 252 minutes. Two stages dominate:
document fetch (118 minutes) and embedding (87 minutes), together about 82\% of
the total. The fetch cost is externally bounded --- EDGAR's fair-access policy
caps automated requests at ten per second per requester
\citep{sec2024edgar} --- so no amount of parallelism removes it.

The filings manifest 1{,}904{,}128 documents in total. We fetched and parsed the
10{,}592 that carry the narrative and exhibit content the benchmark needs,
amounting to 19.7\,GB of raw HTML, which reduces to a 3{,}803\,MiB relational
store plus a 1.0\,GB vector index. \Cref{tab:corpus} gives the resulting counts.

\subsection{Schema}

The relational store is PostgreSQL. Filings and companies are keyed by accession
number and CIK. \texttt{filing\_section} holds one row per extracted Item with its
text; \texttt{section\_chunk} holds the embedded chunks; \texttt{subsidiary} holds
parsed EX-21 rows; \texttt{reporting\_owner} holds parsed ownership rows; and
\texttt{filing\_auditor} holds auditor attestations.

The graph is a projection of the same data, not a second source of truth. It holds
430{,}257 nodes across five labels --- \texttt{Filing} (167{,}050),
\texttt{Subsidiary} (200{,}869), \texttt{Person} (51{,}558), \texttt{Company}
(10{,}511) and \texttt{AuditFirm} (269) --- and 561{,}144 edges across seven
types: \texttt{:SUBSIDIARY\_OF} (212{,}124), \texttt{:FILED} (167{,}050),
\texttt{:OFFICER\_OF} (86{,}180), \texttt{:DIRECTOR\_OF} (70{,}167),
\texttt{:OWNS\_SHARES} (19{,}542), \texttt{:AUDITED\_BY} (5{,}835) and
\texttt{:RESOLVES\_TO} (246).

Two design decisions in that schema carry most of the paper's weight.

\emph{Every edge carries the accession number of the filing that asserts it.}
Provenance is a property of the relationship rather than of a document, which is
what lets an answer cite the exact filing supporting each individual claim rather
than the set of documents that were retrieved.

\emph{\texttt{:RESOLVES\_TO} is the entity-resolution layer, and it is thin.}
It links a subsidiary named in an EX-21 to a company that files in its own right,
a fact no single document states. There are 246 such edges against 200{,}869
subsidiary nodes --- a coverage of 0.12\% --- because EX-21 carries no CIK for a
subsidiary, so the link must be established by name matching. We show in
\Cref{sec:failures} that this thinness is a direct cause of the graph losing the
one category built to exercise it.


\begin{table}[htbp]
\centering
\begin{tabular}{@{}lr@{}}
\toprule
\textbf{Entity} & \textbf{Count} \\
\midrule
Filings (all form types)          & 167,050 \\
Distinct filers (companies)       & 10,511 \\
Documents in filing manifests     & 1,904,128 \\
Documents fetched and parsed      & 10,592 \\
10-K narrative sections extracted & 23,649 \\
Item~1A chunks embedded           & 508,714 \\
EX-21 subsidiary rows             & 212,124 \\
Form 3/4/5 ownership rows         & 150,917 \\
Distinct insiders                 & 51,558 \\
Auditor attestations              & 5,835 \\
\midrule
Graph nodes                       & 430,257 \\
Graph edges                       & 561,144 \\
\midrule
Relational store on disk          & 3,803\,MiB \\
HNSW index                        & 1.0\,GB \\
\bottomrule
\end{tabular}

\caption{The loaded corpus. Twelve monthly EDGAR feeds, 2025-09 through 2026-08.
Node and edge counts are live counts from the graph; the per-type edge breakdown
in \Cref{sec:dataset} sums exactly to the total shown here.}
\label{tab:corpus}
\end{table}

\begin{widetab}[htbp]
\centering
\small
\begin{tabular}{@{}llrrrr@{}}
\toprule
\textbf{Item} & \textbf{Section} & \textbf{Filings} & \textbf{Total chars} & \textbf{Mean} & \textbf{Max} \\
\midrule
1 & Business & 5,539 & 352,550,560 & 63,649 & 717,422 \\
\textbf{1A} & Risk Factors & 5,210 & 608,417,255 & 116,779 & 1,060,461 \\
1B & Unresolved Staff Comments & 63 & 2,000,524 & 31,754 & 127,289 \\
2 & Properties & 2,683 & 12,281,341 & 4,577 & 390,319 \\
3 & Legal Proceedings & 1,518 & 5,126,849 & 3,377 & 117,185 \\
7 & Management's Discussion \& Analysis & 5,609 & 294,551,458 & 52,514 & 473,035 \\
7A & Market Risk & 3,027 & 16,460,821 & 5,438 & 276,843 \\
\bottomrule
\end{tabular}

\caption{Narrative sections extracted per 10-K Item. Only Item~1A (bold) is
embedded, giving 508{,}714 chunks. Its mean length exceeds the encoder's
256-word-piece window by roughly two orders of magnitude, which is why chunking is
forced rather than optional.}
\label{tab:items}
\end{widetab}

\begin{widefig}[htbp]
\centering
\begin{subfigure}[t]{0.48\linewidth}
\centering
\begin{tikzpicture}
\begin{axis}[
  ybar, bar width=15pt, width=\linewidth, height=4.4cm,
  symbolic x coords={{1},{2-5},{6-20},{21-100},{101-500},{500+}}, xtick=data,
  ymin=0, ylabel={Filers}, ylabel near ticks,
  xlabel={Subsidiaries disclosed in EX-21}, xlabel near ticks,
  axis lines*=left, tick style={draw=none}, ymajorgrids, grid style={gray!20},
  nodes near coords, nodes near coords style={font=\tiny}, enlarge x limits=0.13,
]
\addplot[fill=graphcol] coordinates {({1},275) ({2-5},950) ({6-20},1261) ({21-100},1079) ({101-500},439) ({500+},62)};
\end{axis}
\end{tikzpicture}

\caption{EX-21 disclosure per filer.}
\label{fig:subdist}
\end{subfigure}\hfill
\begin{subfigure}[t]{0.48\linewidth}
\centering
\begin{tikzpicture}
\begin{axis}[
  ybar, bar width=7pt, width=\linewidth, height=4.2cm,
  symbolic x coords={{25-08},{25-09},{25-10},{25-11},{25-12},{26-01},{26-02},{26-03},{26-04},{26-05},{26-06},{26-07},{26-08}}, xtick=data,
  x tick label style={rotate=45, anchor=east, font=\tiny},
  ymin=0, ylabel={Filings with parsed documents}, ylabel near ticks,
  axis lines*=left, tick style={draw=none}, ymajorgrids, grid style={gray!20},
  enlarge x limits=0.05,
]
\addplot[fill=veccol] coordinates {({25-08},362) ({25-09},6581) ({25-10},8000) ({25-11},11776) ({25-12},11504) ({26-01},11513) ({26-02},15902) ({26-03},17934) ({26-04},16633) ({26-05},19594) ({26-06},15082) ({26-07},14050) ({26-08},18119)};
\end{axis}
\end{tikzpicture}

\caption{Filings per month.}
\label{fig:months}
\end{subfigure}
\caption{Corpus shape. (a) EX-21 disclosure is extremely skewed --- 275 filers list one subsidiary, 62 list more than 500 --- so any generated \texttt{LIMIT} over a join to this table starves whole entities unless rows are balanced per entity first. (b) The small 2025-08 bucket is spillover, since a monthly feed contains filings dated at the end of the preceding month.}
\label{fig:corpusdist}
\end{widefig}

\begin{figure}[htbp]
\centering
\begin{tikzpicture}
\begin{axis}[
  xbar, bar width=7pt, width=0.74\linewidth, height=6.0cm,
  symbolic y coords={{HNSW + graph},{embed 508,714 chunks},{load RDS},{upload to S3},{extract HTML},{Form 3/4/5},{fetch 10,592 docs},{parse 12 feeds},{fetch feeds}}, ytick=data,
  xmin=0, xmax=132, xlabel={Wall-clock minutes}, xlabel near ticks,
  axis lines*=left, tick style={draw=none}, xmajorgrids, grid style={gray!20},
  nodes near coords, nodes near coords style={font=\tiny}, enlarge y limits=0.09,
  y tick label style={font=\footnotesize},
]
\addplot[fill=graphcol] coordinates {(4.00,{HNSW + graph}) (87.00,{embed 508,714 chunks}) (6.00,{load RDS}) (12.00,{upload to S3}) (19.00,{extract HTML}) (0.40,{Form 3/4/5}) (118.00,{fetch 10,592 docs}) (4.20,{parse 12 feeds}) (0.90,{fetch feeds})};
\end{axis}
\end{tikzpicture}

\caption{Ingestion wall-clock, about 252 minutes total. Document fetch and
embedding are roughly 82\% of it. The fetch component is bounded externally by
EDGAR's ten-requests-per-second fair-access policy and cannot be parallelized
away.}
\label{fig:pipeline}
\end{figure}

% ════════════════════════════════════════════════════════════════════════
\section{System Architecture}
\label{sec:arch}

Three retrieval arms answer the same question over the same stores.
\Cref{fig:arch} shows all three.

\subsection{Text-to-SQL}

The model is shown the complete relational schema, including the
\texttt{subsidiary} and \texttt{reporting\_owner} tables, and writes a SQL query;
the query is executed; the resulting rows are given back to the model, which
writes the answer. Two LLM calls.

This is deliberately a strong baseline rather than a straw man. Because the schema
is complete, every multi-hop answer in the benchmark \emph{is} reachable by joins.
What the model must do unaided is choose the join path, decode a multi-valued
relationship string, and recognize that traversing an ownership chain of unknown
depth requires a recursive CTE. The schema documentation also states the
surname-first convention used for person names --- an asymmetry in an earlier
iteration, where only the graph documentation mentioned it, is recorded among the
harness defects in \Cref{sec:repro}.

\subsection{Vector RAG}

The question is embedded and the $k{=}8$ nearest chunks are retrieved by cosine
similarity from an HNSW index \citep{malkov2018hnsw,pgvector2024}; the chunks are
given to the model, which answers. One LLM call, and no query language, therefore
no repair path.

The encoder is \texttt{all-MiniLM-L6-v2}, a six-layer 22.7M-parameter model
distilled in the MiniLM family \citep{wang2020minilm} and trained as a sentence
encoder \citep{reimers2019sbert}. It emits 384 dimensions with mean pooling and
$L^2$ normalization, so cosine similarity reduces to a dot product.

One configuration detail is load-bearing. The encoder accepts 256 word-pieces,
roughly a thousand characters, while Item 1A averages 116{,}779 characters. The
section therefore cannot be embedded whole: text is chunked at 1{,}400 characters
with 200 characters of overlap, so a single risk-factor section becomes on the
order of a hundred vectors, and 5{,}210 sections become 508{,}714 chunks. We
embed only Item 1A. That is not a handicap imposed on the baseline but the
opposite --- it is the only section any benchmark question asks about, so
restricting the index concentrates the candidate pool and makes the vector arm
\emph{stronger} than it would be with all seven Items embedded.

\subsection{GraphRAG}

Two LLM calls, as with text-to-SQL, but the work between them is split
differently.

\begin{enumerate}
  \item \textbf{Identify.} The model writes a Cypher query whose job is to
        determine \emph{which entities} the question concerns --- typically
        resolving a person to a CIK and finding the companies where they hold a
        role, with their relevant filings.
  \item \textbf{Expand (deterministic).} The system, not the model, expands each
        identified entity: for every company returned, fetch its subsidiaries and
        its risk-factor filings.
  \item \textbf{Fetch text by key.} The accession numbers from the traversal are
        used to select section text by primary key, and windows of $\pm 900$
        characters around matched terms are extracted, at most three per filing.
        This is a lookup, not a search: the graph has already decided which
        documents are in scope, so there is nothing to rank and nothing to miss.
  \item \textbf{Answer.} The model receives the graph facts and the exact
        passages, and writes the answer with citations.
\end{enumerate}

Stage 2 being system-authored rather than LLM-authored is the design's central
claim, and it is a response to a measured failure rather than an aesthetic
preference. An LLM writing a single query over a fan-out join will attach a
\texttt{LIMIT}, and one entity's hundreds of child rows then consume the entire
budget --- after which the model reports, fluently and wrongly, that the other
entities have no subsidiaries or do not exist. Splitting identification from
expansion makes that class of error structurally impossible rather than something
to be prompted against.

\subsection{Shared guardrails}

The same failure threatens both query languages, so the same two transforms are
applied to both, identically. Any trailing \texttt{LIMIT} in a generated query is
raised to 2{,}000 rows before execution; the result set is then balanced to at
most 6 rows per entity and 48 rows overall, where the entity column is detected by
name. Applying these to SQL and Cypher alike means the comparison measures
retrieval architecture rather than which prompt happened to dodge the trap.

Both query-writing arms also get exactly one repair attempt, triggered when a
query raises an exception or returns zero rows: the model sees the error text and
rewrites once. There is no second repair. \Cref{sec:results} shows the two arms
used this allowance very unequally.

\subsection{Model access}

All three arms call a single OpenAI-compatible endpoint provided by Bifrost, an
open-source LLM gateway, which routes to DeepSeek or to AWS Bedrock. The
benchmarked configuration is \texttt{deepseek-v4-flash} at temperature 0 with
reasoning disabled, 700 output tokens for query generation and 1{,}500 for
answers.

Routing through one endpoint is a methodological device, not a convenience. It
makes ``all three arms used the same model with the same settings'' a structural
property of the deployment rather than a convention maintained by discipline: the
model is a request parameter read from one place by all three arms. It also makes
the cross-model probe in \Cref{sec:failures} a configuration change rather than a
code change. Reasoning was disabled because an uncapped reasoning trace consumed
roughly 33{,}000 completion tokens per call, which \texttt{max\_tokens} does not
bound; the setting is identical for every arm, so the comparison is unaffected.


\begin{widefig}[htbp]
\centering
\resizebox{\linewidth}{!}{\begin{tikzpicture}[
  font=\footnotesize,
  >={Stealth[length=1.7mm,width=1.3mm]},
  box/.style={draw=gray!55, line width=0.4pt, rounded corners=1.5pt, fill=white,
              align=center, font=\scriptsize, inner sep=3pt,
              text width=3.9cm, minimum width=4.3cm, minimum height=0.95cm},
  slim/.style={box, minimum height=0.6cm, fill=gray!6, draw=gray!45},
  shared/.style={box, text width=11cm, minimum width=11.5cm, minimum height=0.85cm,
                 fill=gray!8, draw=gray!55},
  gapbox/.style={box, dashed, draw=gray!55, fill=gray!4, text=gray,
                 minimum height=2.55cm},
  key/.style={box, line width=1pt, draw=okcol!70!black, fill=okcol!18,
              minimum height=1.15cm},
  tag/.style={draw=okcol!70!black, fill=okcol!22, rounded corners=1pt,
              inner sep=1.5pt, font=\tiny\bfseries},
  dbicon/.style={cylinder, shape border rotate=90, aspect=0.32, draw=gray!60,
                 fill=gray!18, line width=0.35pt, inner sep=0pt,
                 minimum width=3.8mm, minimum height=2.6mm},
  lane/.style={rounded corners=3pt, inner sep=4pt, line width=0.5pt},
  arr/.style={->, draw=gray!65, line width=0.5pt},
  wire/.style={draw=gray!65, line width=0.5pt},
]

%% ---------- shared question -------------------------------------------------
\node[shared] (q) at (5.05,0)
  {\textbf{Natural-language question over the 10-K / insider corpus}\\[1pt]
   \tiny identical prompt, identical model and $T{=}0$ for all three arms};

%% ---------- lane 1: text-to-SQL --------------------------------------------
\node[box] (l1a) at (0,-2.15)
  {LLM writes SQL\\[1pt]\tiny complete schema in prompt; one repair retry};
\node[box] (l1b) at (0,-3.50)
  {PostgreSQL 17\\[1pt]\tiny execute; 45\,s statement timeout};
\node[slim] (l1c) at (0,-4.70) {result rows \tiny(LIMIT capped at 2{,}000)};
\node[gapbox] (l1d) at (0,-6.70)
  {\itshape no deterministic stage\\[3pt]
   \tiny\upshape whatever join the model happened to write \emph{is} the retrieval:
   one bad predicate silently returns a well-formed, incomplete row set};
\node[box] (l1e) at (0,-8.85)
  {LLM answers\\[1pt]\tiny from the returned rows only};

%% ---------- lane 2: vector RAG ---------------------------------------------
\node[box] (l2a) at (5.05,-2.15)
  {MiniLM-L6-v2 encoder\\[1pt]\tiny 384-d, mean-pooled, normalised (no LLM)};
\node[box] (l2b) at (5.05,-3.50)
  {HNSW cosine top-8\\[1pt]\tiny pgvector over 508{,}714 chunks};
\node[slim] (l2c) at (5.05,-4.70) {8 chunks \tiny(1{,}400 chars, 200 overlap)};
\node[gapbox] (l2d) at (5.05,-6.70)
  {\itshape no query language\\[3pt]
   \tiny\upshape nothing to repair and nothing to expand: the eight nearest
   neighbours are the entire evidence set, whatever the question asked for};
\node[box] (l2e) at (5.05,-8.85)
  {LLM answers\\[1pt]\tiny from the retrieved chunks only};

%% ---------- lane 3: GraphRAG -----------------------------------------------
\node[box] (l3a) at (10.1,-2.15)
  {LLM writes stage-1 Cypher\\[1pt]\tiny \emph{identify entities only}; one repair retry};
\node[box] (l3b) at (10.1,-3.50)
  {Neo4j 5.26 + APOC\\[1pt]\tiny match entity nodes};
\node[slim] (l3c) at (10.1,-4.70) {entity nodes \tiny(seeds, no text yet)};
\node[key] (l3d) at (10.1,-6.00)
  {\textbf{Deterministic expansion}\\[1pt]
   \tiny system code, not the LLM: fixed traversal per seed, 6 filings each, 48 total};
\node[box, minimum height=1cm] (l3e) at (10.1,-7.475)
  {Exact text fetch\\[1pt]\tiny by \texttt{accession\_number} primary key from
   PostgreSQL; 900-char window};
\node[box] (l3f) at (10.1,-8.85)
  {LLM answers\\[1pt]\tiny graph facts + exact snippets};

%% ---------- lane backgrounds ------------------------------------------------
\begin{scope}[on background layer]
  \node[lane, fill=sqlcol!7,   draw=sqlcol!40,   fit=(l1a)(l1b)(l1c)(l1d)(l1e)] (L1) {};
  \node[lane, fill=veccol!7,   draw=veccol!40,   fit=(l2a)(l2b)(l2c)(l2d)(l2e)] (L2) {};
  \node[lane, fill=graphcol!7, draw=graphcol!40, fit=(l3a)(l3b)(l3c)(l3d)(l3e)(l3f)] (L3) {};
\end{scope}

\node[anchor=south west, font=\footnotesize\bfseries, text=sqlcol!75!black]   at (L1.north west) {Text-to-SQL};
\node[anchor=south west, font=\footnotesize\bfseries, text=veccol!75!black]   at (L2.north west) {Vector RAG};
\node[anchor=south west, font=\footnotesize\bfseries, text=graphcol!75!black] at (L3.north west) {GraphRAG};
\node[anchor=south east, font=\tiny, text=gray!65] at (L1.north east) {2 LLM calls};
\node[anchor=south east, font=\tiny, text=gray!65] at (L2.north east) {1 LLM call};
\node[anchor=south east, font=\tiny, text=gray!65] at (L3.north east) {2 LLM calls};

%% ---------- store glyphs ----------------------------------------------------
\node[dbicon, anchor=north east] at ([shift={(-2pt,-2pt)}]l1b.north east) {};
\node[dbicon, anchor=north east] at ([shift={(-2pt,-2pt)}]l2b.north east) {};
\node[dbicon, anchor=north east] at ([shift={(-2pt,-2pt)}]l3b.north east) {};
\node[dbicon, anchor=north east] at ([shift={(-2pt,-2pt)}]l3e.north east) {};

%% ---------- the claim tag ---------------------------------------------------
\node[tag, anchor=south east] at ([yshift=1.5pt]l3d.north east) {not LLM-written};

%% ---------- fan-out from the question --------------------------------------
\coordinate (b1) at (5.05,-0.95);
\draw[wire] (q.south) -- (b1);
\draw[arr] (b1) -| (l1a.north);
\draw[arr] (b1) -| (l2a.north);
\draw[arr] (b1) -| (l3a.north);

%% ---------- in-lane flow ----------------------------------------------------
\draw[arr] (l1a) -- (l1b);  \draw[arr] (l1b) -- (l1c);
\draw[arr] (l1c) -- (l1d);  \draw[arr] (l1d) -- (l1e);
\draw[arr] (l2a) -- (l2b);  \draw[arr] (l2b) -- (l2c);
\draw[arr] (l2c) -- (l2d);  \draw[arr] (l2d) -- (l2e);
\draw[arr] (l3a) -- (l3b);  \draw[arr] (l3b) -- (l3c);
\draw[arr] (l3c) -- (l3d);  \draw[arr] (l3d) -- (l3e);
\draw[arr] (l3e) -- (l3f);

%% ---------- fan-in to the answer -------------------------------------------
\coordinate (b2) at (5.05,-9.80);
\draw[wire] (l1e.south) |- (b2);
\draw[wire] (l2e.south) |- (b2);
\draw[wire] (l3f.south) |- (b2);
\node[shared] (ans) at (5.05,-10.45)
  {\textbf{Answer}\\[1pt]\tiny scored programmatically -- entity recall, required
   keywords, forbidden entities, citation validity (no LLM judge)};
\draw[arr] (b2) -- (ans.north);

\end{tikzpicture}
}
\caption{The three retrieval arms over shared stores. Text-to-SQL and GraphRAG use
two LLM calls each; vector RAG uses one and has no query language to repair. In
GraphRAG only the first hop is LLM-authored: entity expansion is deterministic and
system-authored, which is what prevents one entity's fan-out from consuming the
row budget and the model then reporting that the other entities do not exist.}
\label{fig:arch}
\end{widefig}

% ════════════════════════════════════════════════════════════════════════
\section{Experimental Design}
\label{sec:design}

\subsection{Four question shapes}

The benchmark has 20 questions in four categories of five, listed in full in
\Cref{app:questions}. Each category carries a pre-registered expectation of which
architecture should win, and two of the four were chosen so that the graph is the
\emph{wrong} tool:

\begin{itemize}
  \item \textbf{T1, single-document semantic.} What did one named company disclose
        about one topic in its risk factors? The whole answer lies inside one
        document. Expected winner: vector RAG or GraphRAG.
  \item \textbf{T2, structured aggregation.} Which audit firm serves the most
        companies; which jurisdiction of incorporation is most common. These are
        \texttt{GROUP BY} queries over typed columns. Expected winner:
        text-to-SQL, and a traversal should add nothing.
  \item \textbf{T3, multi-hop relational plus text.} For every company where a
        named person holds an insider role, list its subsidiaries and summarize
        its reported cybersecurity risks. Expected winner: GraphRAG.
  \item \textbf{T4, cross-document entity resolution.} Is a named subsidiary also
        an SEC filer in its own right, and which parent lists it? Expected
        winner: GraphRAG.
\end{itemize}

\subsection{Selecting subjects that make identity hard}

The T3 subjects were not chosen for convenience. Two conditions had to hold
simultaneously, both checked against the full-text index over all 5{,}210
risk-factor sections:

\begin{enumerate}
  \item the surname appears in \emph{zero} risk-factor sections, so no similarity
        search can reach the person at all; and
  \item two or more distinct filer identifiers share that surname, so resolving
        identity by name is demonstrably wrong rather than merely imprecise.
\end{enumerate}

Candidates were rejected on the first condition when the surname was too common in
the text (32, 49, 63 and 109 matching sections respectively) and on the second
when the surname turned out to be unique in the corpus. One unique-surname subject
was retained deliberately as a control: same question shape, no identity trap.

\subsection{Scoring}

Scoring is a deterministic function of the answer text. There is no LLM judge, for
the reason given in \Cref{sec:related}. An attempt \textbf{passes} when all four
conditions hold:

\begin{enumerate}
  \item every required entity is named (entity recall $=1.0$);
  \item at least one required topical keyword is present;
  \item no forbidden entity is named --- these are the specific wrong companies
        that a name-based rather than identifier-based resolution would return;
  \item where ground-truth citations exist, at least one is correct.
\end{enumerate}

The fourth condition is what makes provenance count. Without it, a model that
named the right companies but cited nothing would score identically to one that
cited the exact filings. Failures are then labelled: \textbf{falsehood} when the
answer is not correct, is not a refusal, and either names a forbidden entity or
cites a filing outside the valid set; \textbf{uncited} when the required entities
and keywords are present with nothing false but no correct citation;
\textbf{refused} when a regex over the answer detects an explicit statement that
the question cannot be answered from the evidence; and \textbf{fail} otherwise.
The labels partition the attempts, so the columns of \Cref{tab:main} sum to 60.

Aggregate questions have no ground-truth citation set. For those, any accession the
model cites is legitimate provenance and is not counted against it --- an earlier
version of the scorer did count it, and \Cref{sec:repro} records what that cost.

\subsection{Three attempts, all of them counted}

Every question was asked three times of every architecture, giving
$20 \times 3 \times 3 = 180$ scored attempts. The three attempts are
\emph{independent repeats to measure consistency, not retries}: every attempt is
scored and every attempt counts, and there is no best-of-three anywhere in the
reported numbers.

This is not ceremony. Even at temperature 0 the generations are not
deterministic, and \textbf{9 of the 60}
question-by-architecture cells returned different labels across their three
attempts --- 7 for text-to-SQL, 2 for GraphRAG and
0 for vector RAG. Text-to-SQL was inconsistent on four of the five
multi-hop questions, which is precisely where its false statements occur: an
evaluation drawing a single attempt per question could record an honest refusal
on the same question where another draw asserts something false.

We therefore report two metrics. \textbf{Attempts passed} out of 60 is the raw
rate. \textbf{Questions solved} out of 20 credits a question only when all three
of its attempts passed, and is the stricter reading: an answer that is correct
once in three is not a capability a practitioner can deploy. The two are not
related by a fixed ratio --- an architecture with $q$ questions solved has at
least $3q$ attempts passed, and the excess is contributed by questions solved
inconsistently.

\subsection{Fairness controls}

A rigged baseline would prove nothing, so we record exactly what was equalized:
identical model, temperature, reasoning setting and answer-token budget for all
three arms; an equal LLM call budget of two calls for the query-writing arms and
one for vector RAG, which has no query to write; the complete schema shown to
text-to-SQL so that every multi-hop answer is reachable by joins; the
surname-first naming convention documented in both schema descriptions; one
repair attempt for each query-writing arm; and the identical fan-out guardrail on
both query languages. Two of the four categories were selected so that the graph
should lose them.

We note one asymmetry we did not eliminate, because it favours the graph and we
would rather report it than hide it. The repair trigger fires on an exception or
zero rows. An over-broad SQL join returns \emph{many} rows and therefore never
trips it, even when the answer is wrong, whereas an exact graph property match
that misses returns nothing and earns a free second attempt. The same rule confers
unequal benefit; \Cref{tab:repair} quantifies the result.


% ════════════════════════════════════════════════════════════════════════
\section{Results}
\label{sec:results}

\Cref{tab:main} gives the headline numbers and \Cref{tab:bytype} the per-category
breakdown. \Cref{fig:bytype} plots the latter and \Cref{tab:grid} shows every one
of the 180 attempts individually, unaggregated.

Summed over the question set, GraphRAG solved 14 of 20 questions on all three
attempts and passed 45 of 60 attempts; text-to-SQL solved 9 and passed 29; vector
RAG solved 3 and passed 9. These totals weight four hand-chosen categories
equally and GraphRAG lost two of the four, so the aggregate is a property of this
question mix rather than an overall ranking of the architectures. The per-category
results are the informative ones.

\subsection{By category}

\textbf{T1, single-document semantic.} GraphRAG passed 15 of 15. Vector RAG, the
expected winner, passed 3 of 15 --- all three from a single question --- with 6
refusals and 6 answers that named the right company and the right risks but cited
no filing. Text-to-SQL passed 0 of 15 despite reaching the correct text in 14
attempts: its generated queries used full-text ranking to find the passage but did
not select the accession number, so the answers were right and unciteable. This is
the clearest demonstration that provenance is a separate capability from
retrieval.

\textbf{T2, structured aggregation.} Text-to-SQL passed 15 of 15, as predicted.
GraphRAG passed 6 of 15, solving 2 of 5 questions, with 7 attempts returning wrong
answers and 2 refusing. Vector RAG passed 0 of 15 and refused all fifteen times,
which is the correct behaviour: an aggregate over 10{,}511 companies is not
present in any eight retrieved chunks, and the arm said so rather than guessing.

\textbf{T3, multi-hop relational plus text.} GraphRAG passed 15 of 15; both
baselines passed 0 of 15. This is the category the design targets and the result
is unambiguous. Text-to-SQL produced the benchmark's only two falsehoods here.

\textbf{T4, cross-document entity resolution.} Written for the graph, and the
graph came second. Text-to-SQL passed 14 of 15 and solved 4 of 5 questions;
GraphRAG passed 9 of 15 and solved 2 of 5 with 2 more solved inconsistently, its
6 misses all refusals rather than errors; vector RAG passed 6 of 15.

\subsection{Outcome composition}

\Cref{fig:outcomes} decomposes all 60 attempts per arm. The composition matters as
much as the pass rate, because the failure modes are not interchangeable.

Vector RAG refused 45 times. That is a low score but a benign failure: asked about
a company whose chunks it could not find, it generally said the excerpts did not
contain the answer rather than answering about whichever company did surface.
Text-to-SQL was uncited 16 times --- right, and unverifiable. GraphRAG was never
uncited, which follows from its construction: the traversal returns accession
numbers before any text is read, so every claim arrives with a key attached.

Only \textbf{2 of the 180 attempts} asserted something verifiably false, both
text-to-SQL on T3.

\subsection{Cost and latency}

\Cref{fig:cost} plots tokens against attempts passed. GraphRAG consumed
780{,}500 tokens to text-to-SQL's 158{,}171 and vector RAG's 118{,}850 --- about
17{,}300 tokens per passing attempt against 5{,}450 and 13{,}200 respectively. Its
accuracy is bought, not free. The extreme case is T2, where GraphRAG spent
590{,}453 tokens against text-to-SQL's 24{,}604: roughly 24 times the cost for
40\% of the accuracy, on the category it was expected to lose.

The reverse cost is larger and easier to overlook. \Cref{fig:latency} plots mean
latency on a log scale. Text-to-SQL averaged 25.3\,s overall against GraphRAG's
5.9\,s and vector RAG's 1.6\,s, but the overall mean hides the shape: on T3 it
averaged 92\,s and its worst single attempt took 445\,s. The cause is not the
model but the SQL --- the generated queries join \texttt{subsidiary},
\texttt{filing\_section} and \texttt{reporting\_owner} without narrowing first.
Unbounded, LLM-authored SQL is a cost risk against one's own warehouse
independently of whether the answer is correct.

\subsection{Repair usage}

\Cref{tab:repair} and \Cref{fig:repair} report how often each arm used its single
repair allowance. GraphRAG triggered a repair on 34 of 60 attempts against
text-to-SQL's 7 --- nearly five times as often --- and on T4 it needed one on
\emph{all fifteen} attempts. Vector RAG has no query to repair and used none.

This asymmetry accounts for much of the token gap, and it should temper the
reading of GraphRAG's wins: writing correct Cypher against this schema was harder
for the model than writing correct SQL, and the graph arm leaned on its second
chance far more heavily. As noted in \Cref{sec:design}, the zero-row trigger is
also not neutral between the two languages, and the direction of that bias favours
the graph.


\begin{widetab}[htbp]
\centering
\begin{tabular}{@{}lccc@{}}
\toprule
& \textbf{Text-to-SQL} & \textbf{Vector RAG} & \textbf{GraphRAG} \\
\midrule
Questions solved on all 3 attempts & 10/20 & 3/20 & \textbf{16/20} \\
Questions solved inconsistently & 3 & 0 & 1 \\
Attempts passed (of 60) & 36 & 9 & \textbf{50} \\
\midrule
Confident falsehoods & \textcolor{badcol}{6} & \textbf{0} & \textbf{0} \\
Correct but uncited & 10 & 12 & \textbf{0} \\
Explicit refusals & 8 & 39 & 1 \\
\midrule
Mean evidence supplied (chars) & 1,795 & 8,557 & 6,084 \\
Mean latency (s) & 10.0 & \textbf{1.3} & 4.6 \\
Max latency (s) & 274.5 & 4.6 & 8.6 \\
Total tokens & 144,068 & \textbf{118,185} & 520,929 \\
Tokens per passing attempt & \textbf{4,001} & 13,131 & 10,418 \\
LLM calls & 125 & 60 & 139 \\
\bottomrule
\end{tabular}

\caption{Aggregate results over 180 scored attempts. \emph{Questions solved}
credits a question only when all three of its attempts passed. Bold marks the best
value in a row. The outcome labels partition the attempts, so pass, falsehood,
uncited, refused and the residual failures sum to 60 per arm.}
\label{tab:main}
\end{widetab}

\begin{widetab}[htbp]
\centering
\small
\begin{tabular}{@{}ll rr rr rr@{}}
\toprule
& & \multicolumn{2}{c}{\textbf{Text-to-SQL}} & \multicolumn{2}{c}{\textbf{Vector RAG}} & \multicolumn{2}{c}{\textbf{GraphRAG}} \\
\cmidrule(lr){3-4}\cmidrule(lr){5-6}\cmidrule(lr){7-8}
\textbf{Type} & \textbf{Question class} & Q & A & Q & A & Q & A \\
\midrule
T1 & single-document semantic & 0/5 & 0/15 & 1/5 & 3/15 & \textbf{5}/5 & 15/15 \\
T2 & structured aggregation & \textbf{5}/5 & 15/15 & 0/5 & 0/15 & 2/5 & 6/15 \\
T3 & multi-hop relational + text & 0/5 & 0/15 & 0/5 & 0/15 & \textbf{5}/5 & 15/15 \\
T4 & cross-document entity resolution & \textbf{4}/5 & 14/15 & 2/5 & 6/15 & 2/5 & 9/15 \\
\bottomrule
\end{tabular}

\caption{Per-category results. \textbf{Q}~=~questions solved on all three
attempts (of 5); \textbf{A}~=~attempts passed (of 15). Two of the four categories
were pre-registered as ones the graph should lose. It lost T2 as expected and T4
against expectation.}
\label{tab:bytype}
\end{widetab}

\begin{widefig}[htbp]
\centering
\begin{subfigure}[t]{0.48\linewidth}
\centering
\begin{tikzpicture}
\begin{axis}[
  ybar, bar width=8pt, width=\linewidth, height=5.2cm,
  symbolic x coords={T1,T2,T3,T4}, xtick=data,
  ymin=0, ymax=16, ytick={0,5,10,15},
  ylabel={Attempts passed (of 15)}, ylabel near ticks,
  legend style={at={(0.5,1.22)}, anchor=north, legend columns=3, draw=none, font=\footnotesize},
  nodes near coords, nodes near coords style={font=\tiny},
  axis lines*=left, tick style={draw=none}, ymajorgrids, grid style={gray!20},
  enlarge x limits=0.15,
]
\addplot[fill=sqlcol] coordinates {(T1,0) (T2,15) (T3,0) (T4,14)};
\addplot[fill=veccol] coordinates {(T1,3) (T2,0) (T3,0) (T4,6)};
\addplot[fill=graphcol] coordinates {(T1,15) (T2,6) (T3,15) (T4,9)};
\legend{Text-to-SQL, Vector RAG, GraphRAG}
\end{axis}
\end{tikzpicture}

\caption{Attempts passed per category.}
\label{fig:bytype}
\end{subfigure}\hfill
\begin{subfigure}[t]{0.48\linewidth}
\centering
\begin{tikzpicture}
\begin{axis}[
  ybar stacked, bar width=26pt, width=\linewidth, height=5.2cm,
  symbolic x coords={{Text-to-SQL},{Vector RAG},{GraphRAG}}, xtick=data,
  ymin=0, ymax=61, ylabel={Attempts (of 60)}, ylabel near ticks,
  legend style={at={(0.5,1.22)}, anchor=north, legend columns=5, draw=none, font=\tiny},
  axis lines*=left, tick style={draw=none}, ymajorgrids, grid style={gray!20},
  enlarge x limits=0.28, x tick label style={font=\footnotesize},
]
\addplot[fill=okcol] coordinates {({Text-to-SQL},29) ({Vector RAG},9) ({GraphRAG},45)};
\addplot[fill=badcol] coordinates {({Text-to-SQL},2) ({Vector RAG},0) ({GraphRAG},0)};
\addplot[fill=warncol] coordinates {({Text-to-SQL},16) ({Vector RAG},6) ({GraphRAG},0)};
\addplot[fill=gray!45] coordinates {({Text-to-SQL},12) ({Vector RAG},45) ({GraphRAG},8)};
\addplot[fill=gray!20] coordinates {({Text-to-SQL},1) ({Vector RAG},0) ({GraphRAG},7)};
\legend{pass, falsehood, uncited, refused, other fail}
\end{axis}
\end{tikzpicture}

\caption{How all 60 attempts resolved.}
\label{fig:outcomes}
\end{subfigure}
\caption{Results. The composition in (b) matters as much as the rate in (a): vector RAG's 45 refusals are a benign failure, text-to-SQL's 16 uncited answers are right but unverifiable, and GraphRAG was never uncited because the traversal returns accession numbers before any text is read.}
\label{fig:results}
\end{widefig}

\begin{widefig}[htbp]
\centering
\begin{subfigure}[t]{0.48\linewidth}
\centering
\begin{tikzpicture}
\begin{semilogyaxis}[
  ybar, bar width=8pt, width=\linewidth, height=4.8cm,
  symbolic x coords={T1,T2,T3,T4}, xtick=data,
  ymin=0.6, ymax=400, ylabel={Mean latency (s, log)}, ylabel near ticks,
  legend style={at={(0.5,1.25)}, anchor=north, legend columns=3, draw=none, font=\footnotesize},
  nodes near coords, nodes near coords style={font=\tiny},
  point meta=explicit symbolic,
  axis lines*=left, tick style={draw=none}, ymajorgrids, grid style={gray!20},
  enlarge x limits=0.15, log origin=infty,
]
\addplot[fill=sqlcol] coordinates {(T1,3.50) [3.5] (T2,1.90) [1.9] (T3,31.60) [31.6] (T4,2.90) [2.9]};
\addplot[fill=veccol] coordinates {(T1,1.30) [1.3] (T2,1.20) [1.2] (T3,1.50) [1.5] (T4,1.40) [1.4]};
\addplot[fill=graphcol] coordinates {(T1,3.90) [3.9] (T2,4.90) [4.9] (T3,5.00) [5.0] (T4,4.80) [4.8]};
\legend{Text-to-SQL, Vector RAG, GraphRAG}
\end{semilogyaxis}
\end{tikzpicture}

\caption{Mean latency per category (log).}
\label{fig:latency}
\end{subfigure}\hfill
\begin{subfigure}[t]{0.48\linewidth}
\centering
\begin{tikzpicture}
\begin{axis}[
  width=\linewidth, height=5.4cm,
  xlabel={Total tokens consumed (thousands)}, ylabel={Attempts passed (of 60)},
  ylabel near ticks, xlabel near ticks,
  xmin=0, xmax=900, ymin=0, ymax=56,
  axis lines*=left, tick style={draw=none}, grid=both, grid style={gray!20},
]
\addplot[only marks, mark=*, mark size=2.9pt, sqlcol] coordinates {(158.2,29)};
\addplot[only marks, mark=*, mark size=2.9pt, veccol] coordinates {(118.8,9)};
\addplot[only marks, mark=*, mark size=2.9pt, graphcol] coordinates {(780.5,45)};
\node[anchor=north west, align=left, font=\scriptsize, inner sep=1pt, xshift=3mm]
  at (axis cs:158.2,29) {\textbf{Text-to-SQL}\\\tiny 5,454 tok/pass};
\node[anchor=south west, align=left, font=\scriptsize, inner sep=1pt, xshift=3mm]
  at (axis cs:118.8,9) {\textbf{Vector RAG}\\\tiny 13,205 tok/pass};
\node[anchor=north east, align=right, font=\scriptsize, inner sep=1pt, xshift=-3mm]
  at (axis cs:780.5,45) {\textbf{GraphRAG}\\\tiny 17,344 tok/pass};
\end{axis}
\end{tikzpicture}

\caption{Tokens against attempts passed.}
\label{fig:cost}
\end{subfigure}
\caption{The economics. (a) Text-to-SQL's T3 mean is the SQL and not the model: the generated queries join three tables without narrowing first. (b) GraphRAG's accuracy is bought: it has the highest token cost per passing attempt of the three, the exact figures being in \Cref{tab:main}.}
\label{fig:econ}
\end{widefig}

\begin{table}[htbp]
\centering
\begin{tabular}{@{}lccccc@{}}
\toprule
\textbf{Architecture} & \textbf{Overall} & \textbf{T1} & \textbf{T2} & \textbf{T3} & \textbf{T4} \\
\midrule
Text-to-SQL & 5/60 & 3/15 & 1/15 & 0/15 & 1/15 \\
Vector RAG & 0/60 & 0/15 & 0/15 & 0/15 & 0/15 \\
GraphRAG & 19/60 & 1/15 & 9/15 & 0/15 & 9/15 \\
\bottomrule
\end{tabular}

\caption{Use of the single repair allowance, which both query-writing arms
receive on identical terms. GraphRAG leaned on it nearly five times as often, and
on T4 needed it on every attempt. Vector RAG has no query language to repair.}
\label{tab:repair}
\end{table}

\begin{table}[htbp]
\centering
\scriptsize
\begin{tabular}{@{}l ccc@{}}
\toprule
\textbf{Q} & \textbf{Text-to-SQL} & \textbf{Vector RAG} & \textbf{GraphRAG} \\
\midrule
\texttt{T1-1} & \gunc\gunc\gunc & \gref\gunc\gunc & \gpass\gpass\gpass \\
\texttt{T1-2} & \gunc\gunc\gunc & \gpass\gpass\gpass & \gpass\gpass\gpass \\
\texttt{T1-3} & \gunc\gunc\gunc & \gref\gunc\gref & \gpass\gpass\gpass \\
\texttt{T1-4} & \gunc\gunc\gunc & \gref\gref\gref & \gpass\gpass\gpass \\
\texttt{T1-5} & \gunc\gunc\gref & \gunc\gunc\gunc & \gpass\gpass\gpass \\
\addlinespace[2pt]
\texttt{T2-1} & \gpass\gpass\gpass & \gref\gref\gref & \gfail\gfail\gfail \\
\texttt{T2-2} & \gpass\gpass\gpass & \gref\gref\gref & \gpass\gpass\gpass \\
\texttt{T2-3} & \gpass\gpass\gpass & \gref\gref\gref & \gref\gfail\gref \\
\texttt{T2-4} & \gpass\gpass\gpass & \gref\gref\gref & \gfail\gfail\gfail \\
\texttt{T2-5} & \gpass\gpass\gpass & \gref\gref\gref & \gpass\gpass\gpass \\
\addlinespace[2pt]
\texttt{T3-1} & \gref\gref\gref & \gref\gref\gref & \gpass\gpass\gpass \\
\texttt{T3-2}\,$\star$ & \ghall\gfail\gref & \gref\gref\gref & \gpass\gpass\gpass \\
\texttt{T3-3} & \ghall\gref\gref & \gref\gref\gref & \gpass\gpass\gpass \\
\texttt{T3-4} & \gref\gref\gref & \gref\gref\gref & \gpass\gpass\gpass \\
\texttt{T3-5} & \gref\gunc\gunc & \gref\gref\gref & \gpass\gpass\gpass \\
\addlinespace[2pt]
\texttt{T4-1} & \gpass\gpass\gpass & \gref\gref\gref & \gpass\gpass\gpass \\
\texttt{T4-2} & \gpass\gpass\gpass & \gref\gref\gref & \gref\gref\gref \\
\texttt{T4-3} & \gpass\gpass\gpass & \gref\gref\gref & \gpass\gref\gpass \\
\texttt{T4-4} & \gref\gpass\gpass & \gpass\gpass\gpass & \gref\gpass\gref \\
\texttt{T4-5} & \gpass\gpass\gpass & \gpass\gpass\gpass & \gpass\gpass\gpass \\
\bottomrule
\end{tabular}

\caption{Every one of the 180 attempts, unaggregated; three glyphs per cell are
the three attempts. \gpass~pass, \ghall~falsehood, \gunc~correct but uncited,
\gref~refused, \gfail~other failure. $\star$ marks the featured multi-hop
question. 10 of these 60 cells did not return the same label on all three
attempts.}
\label{tab:grid}
\end{table}

\begin{figure}[htbp]
\centering
\begin{tikzpicture}
\begin{axis}[
  ybar, bar width=10pt, width=\linewidth, height=4.6cm,
  symbolic x coords={T1,T2,T3,T4}, xtick=data,
  ymin=0, ymax=16, ytick={0,5,10,15},
  ylabel={Attempts needing repair (of 15)}, ylabel near ticks,
  legend style={at={(0.5,1.25)}, anchor=north, legend columns=2, draw=none, font=\footnotesize},
  nodes near coords, nodes near coords style={font=\tiny},
  axis lines*=left, tick style={draw=none}, ymajorgrids, grid style={gray!20},
  enlarge x limits=0.18,
]
\addplot[fill=sqlcol] coordinates {(T1,3) (T2,1) (T3,0) (T4,1)};
\addplot[fill=graphcol] coordinates {(T1,1) (T2,9) (T3,0) (T4,9)};
\legend{Text-to-SQL, GraphRAG}
\end{axis}
\end{tikzpicture}

\caption{Repair frequency by category. GraphRAG required a repair on all fifteen
T4 attempts, which is the entity-resolution category it lost.}
\label{fig:repair}
\end{figure}

% ════════════════════════════════════════════════════════════════════════
\section{Failure Analysis}
\label{sec:failures}

The failures are the most informative part of this study, and two of them
contradict what we expected before running it.

\subsection{Vector RAG lost the category it was expected to win}

T1 asks what one named company disclosed about one topic in its risk factors. The
answer is entirely inside a single document, retrieval is a ranking problem, and
similarity search should be at its best. It passed 3 of 15 attempts, all from one
question.

The mechanism is visible in the retrieved set. Asked what manufacturing and
supply-chain risks \emph{Moderna} disclosed, the top-8 chunks came from Entegris,
Hayward Holdings, Open Text, Arrowhead Pharmaceuticals and NVIDIA. Every one of
those chunks genuinely discusses supply-chain risk. The retriever did what it was
asked. The problem is that ``belongs to Moderna'' is not a property the embedding
encodes, so the company name in the question cannot outweigh topical similarity
across 508{,}714 chunks drawn from 5{,}210 filers.

The decisive detail is that \textbf{this question passed on a single month of
filings}. The failure appeared when the corpus grew. It is therefore a
scale-induced limitation of entity filtering by similarity, not a tuning error, and
it is the strongest argument in this paper for structural retrieval --- and one we
did not anticipate.

We are careful about the boundary of this claim. This experiment cannot fully
separate ``embeddings cannot filter by entity'' from ``this retriever
configuration cannot'', because the 1{,}400-character chunks exceed the encoder's
256-token window and $k$ was fixed at 8. A hybrid retriever with a metadata filter
on CIK would very likely fix T1 --- which is precisely the point, since a metadata
filter on a stable identifier \emph{is} structural retrieval, arriving by another
route.

\subsection{Capability and safety moved in opposite directions}

The benchmark's most uncomfortable result is not a score but a direction of
travel. While calibrating the harness we found that both stores normalise names
by stripping corporate suffixes and that neither schema description said so
(\Cref{sec:harness}). Documenting it identically for every arm --- one sentence
of schema prose, no change to retrieval code, questions or scoring --- moved
three of the four categories.

It also made the strongest baseline \emph{more dangerous}. With the schema fully
described, text-to-SQL stopped refusing multi-hop questions and started answering
them, and its confident falsehoods rose accordingly to the six reported here.
Better schema knowledge did not make it more careful. It converted honest
abstentions into assertive answers, and on this corpus an assertive multi-hop
answer keyed on a person's name is exactly how three different people called
Lozano become one.

This is worth separating from the accuracy comparison, because the two do not
move together. Giving a text-to-SQL system a better description of its own schema
is an unambiguous capability improvement and it made the system less safe on the
questions where identity is ambiguous. GraphRAG's falsehood count was zero
throughout, not because it is a better writer of queries --- it needed its repair
allowance far more often --- but because identity is resolved on an identifier
before any text is read, so the wrong people are excluded structurally rather
than by the model choosing correctly.

\subsection{GraphRAG lost aggregation, as predicted}

Text-to-SQL scored 15 of 15 on T2 and GraphRAG 6 of 15. Unlike the previous two,
this outcome was expected --- T2 exists in the benchmark precisely because a
traversal should add nothing to a \texttt{GROUP BY} --- but the manner of the loss
is instructive.

Asked for the most common jurisdiction of incorporation among all disclosed
subsidiaries, the generated Cypher grouped by company rather than by jurisdiction,
returned per-company counts, and the model reported Delaware with 2{,}415
subsidiaries. Delaware is the right answer; 2{,}415 is one company's count rather
than the 62{,}504 across the corpus. Right entity, wrong aggregate. The answer
also attached 190 accession numbers to a question that needs none.

For counting over typed columns the graph adds no information and the extra hop
adds a place to go wrong, at 24 times the token cost.

\subsection{Every falsehood was a surname collision}

Six of the 180 attempts stated something verifiably false. All six were
text-to-SQL, all six on the multi-hop category, and all six are the same error,
spread across four of its five questions (\texttt{T3-1}, \texttt{T3-2}, \texttt{T3-3}, \texttt{T3-4}). Asked about \emph{Monica Lozano},
the model reported companies including Mondelez and Ternium; those belong to
Mariano and Santiago Lozano, different people who share a surname. Asked about
\emph{William Giles}, it reported ten companies.

The SQL was not malformed and the model was fluent and confident. It matched on
name because name is what the question supplied and what the text contains. This
is the exact failure that resolving identity to an identifier before reading any
text prevents structurally rather than probabilistically: GraphRAG passed 14 of
15 attempts on this category and made no false statement anywhere in the
benchmark, because the other Lozanos never entered its result set.

One caution about single-attempt evaluation follows. Of the three attempts
text-to-SQL made on the featured question the outcomes were not identical, and an
evaluation that drew only one of them could easily have recorded an honest
refusal and missed the failure entirely. Four of the five multi-hop questions
produced at least one falsehood and at least one non-falsehood, so on this
category a single draw per question is close to a coin toss over whether the most
important failure mode in the paper is observed at all.

\subsection{The verdict is model-sensitive}

Because all arms route through one gateway (\Cref{sec:arch}), we probed the same
GraphRAG pipeline on other models. \Cref{tab:crossmodel} reports the featured
multi-hop question: the benchmarked model passes, Claude Haiku 4.5 fails because
its stage-1 Cypher collapsed to a single company and missed two others, and Amazon
Nova 2 Lite refuses.

This is a probe and not a measurement --- one question, one attempt each, no
re-scoring of the full set --- and we draw only the negative conclusion it
supports: a single-model benchmark bounds what can be claimed. The architecture
constrains what evidence reaches the model, but the model still has to write a
correct traversal, and not all of them do.


\begin{widetab}[htbp]
\centering
\small
\begin{tabular}{@{}llp{0.44\linewidth}@{}}
\toprule
\textbf{Model} & \textbf{Verdict} & \textbf{Observed behaviour} \\
\midrule
DeepSeek V4 Flash (benchmarked) & pass    & correct three-company traversal \\
Claude Haiku 4.5 (Bedrock)      & fail    & stage-1 Cypher returned \textsc{target corp} alone; Apple and Bank of America missing \\
Amazon Nova 2 Lite (Bedrock)    & refused & produced no usable traversal \\
\bottomrule
\end{tabular}

\caption{Cross-model probe on the featured multi-hop question with the GraphRAG
pipeline held constant. One question and one attempt per model, so this is a probe
rather than a measurement; it supports only the negative conclusion that a
single-model benchmark bounds what can be claimed.}
\label{tab:crossmodel}
\end{widetab}

% ════════════════════════════════════════════════════════════════════════
\section{Threats to Validity}
\label{sec:threats}

\textbf{One corpus, one domain.} Every measurement comes from a single snapshot of
one filing system. SEC disclosure is unusually favourable to structural retrieval:
it is mandated, schematized, and carries a stable public entity identifier in the
CIK. Corpora without an authoritative identifier --- internal wikis, support
tickets, email --- would require the entity-resolution step this paper largely
assumes, and \Cref{sec:failures} shows that step is our weakest layer. Nothing
here should be read as generalizing to such corpora.

\textbf{Twenty questions is a small sample.} Each category rests on five
questions, so a category-level rate has a wide confidence interval: a single
question changes a category by 20 percentage points at question level and 3
attempts out of 15 at attempt level. The category results are directional, not
conclusive. The one result we would defend as strong is T3, where the margin is
15 of 15 against 0 of 15 for both baselines and the mechanism is independently
diagnosed.

\textbf{The questions were written by the author of the graph.} This is the most
serious threat and it cannot be fully mitigated. We took three measures: two of
the four categories were pre-registered as ones the graph should lose, and it did
lose both; the T3 subjects were selected by a measured criterion rather than by
inspection of results (\Cref{sec:design}); and the ground truth was verified
against the loaded tables rather than against model outputs. None of this
eliminates the risk that the question set favours the design in ways we cannot
see. An independent question set is the obvious next step.

\textbf{One model, and the verdict moves with it.} All reported numbers come from
\texttt{deepseek-v4-flash}. The probe in \Cref{tab:crossmodel} shows the featured
question passing on one model, failing on another and being refused by a third
with the retrieval architecture held constant. Our results characterize
\emph{architecture $\times$ model}, and we did not vary the model across the full
180 attempts.

\textbf{The graph inherits EX-21 incompleteness.} Rule 601(b)(21)(ii) allows
filers to omit insignificant subsidiaries, so \texttt{:SUBSIDIARY\_OF} is a floor
on corporate structure. Any question whose correct answer depends on a subsidiary
the filer chose not to disclose is unanswerable from this corpus by any of the
three arms, and our ground truth is defined against what was disclosed rather than
against reality.

\textbf{Category results move with implementation details.} While calibrating
the harness we found an undocumented name-normalisation convention in the schema
descriptions; documenting it changed three of the four categories
(\Cref{sec:harness}). Nothing about any retrieval architecture changed. If one
sentence of schema prose was worth three categories, other such sensitivities are
plausibly latent in results we have not thought to re-examine, and the ones we
would be least likely to notice are those that flatter the design. The same logic
applies to the wins: some part of GraphRAG's margins may reflect implementation
effort we spent on the graph and not on the baselines.
\texttt{:RESOLVES\_TO} coverage of 0.12\% remains a hard bound on what any
entity-resolution question can measure here.

\textbf{Repair asymmetry favours the graph.} As recorded in \Cref{sec:design},
the zero-row repair trigger fires for a graph query that misses but not for an
over-broad SQL join that returns many wrong rows. GraphRAG used its repair on 34
of 60 attempts against text-to-SQL's 7, so it received materially more effective
assistance from an identical rule.

\textbf{A single year, and a moving corpus.} The window spans twelve monthly
feeds. Annual filings cluster seasonally (\Cref{fig:months}), so the mix of form
types is not stationary within it, and the absolute counts are not reproducible
against a later snapshot because EDGAR accumulates.

\textbf{Undetected errors are likely.} We found several defects in our own
harness, one of them a ground-truth error discovered after we believed the
results were final (\Cref{sec:harness}). That we found it late is direct evidence
that others may remain, and we have no way to bound how many.


% ════════════════════════════════════════════════════════════════════════
\section{Reproducibility and Engineering Cost}
\label{sec:repro}

\subsection{Software and infrastructure}

Both relational and vector retrieval run on PostgreSQL 17 with pgvector 0.8.2
\citep{pgvector2024}, using an HNSW index over \texttt{vector\_cosine\_ops}
\citep{malkov2018hnsw} alongside \texttt{pg\_trgm} and a GIN \texttt{tsvector}
index. \texttt{hnsw.iterative\_scan} is set to \texttt{relaxed\_order} for reasons
given below. The graph is Neo4j 5.26 Community with APOC \citep{neo4j2024}.
Embeddings come from \texttt{all-MiniLM-L6-v2} at 384 dimensions with a 256
word-piece limit, chunked at 1{,}400 characters with 200 of overlap, batch size
256. Generated SQL and Cypher execute under a 45-second statement timeout.

The deployment is a single application host --- an AWS \texttt{m7g.xlarge}, 4 vCPU
and 16\,GB, running the graph, the API and the gateway --- with PostgreSQL on a
\texttt{db.t4g.medium} instance and CloudFront terminating TLS. There is no NAT
gateway; object-store access goes through a VPC endpoint. \Cref{tab:cost} gives
the monthly cost at roughly \$165. The embedding pass was run on a temporarily
resized 16-vCPU host at 97 chunks/s and the instance was then scaled back.

Ingestion takes about 252 minutes end to end (\Cref{fig:pipeline}), of which
document fetch and embedding are about 82\%. The fetch component is bounded
externally by EDGAR's ten-requests-per-second fair-access policy
\citep{sec2024edgar} and cannot be parallelized away.

\subsection{Auditing the harness}
\label{sec:harness}

Building the measurement code surfaced defects that, unnoticed, would each have
produced a publishable-looking but incorrect result: a pgvector post-filter that
returned 0 of 8 rows while 164 matched, so the vector arm was briefly scored on
an empty retriever; an answer-token cap that truncated answers mid-citation,
silently converting passes into uncited; a global token counter that attributed
every thread's cost to whichever arm read it first. These were found and fixed
before the reported numbers were generated, so they are development history
rather than caveats on the results, and we do not enumerate them here.

Two are worth reporting, because they did reach published numbers.

\textbf{A ground-truth error, found after the first analysis.} One
cross-document question accepted a single expected parent where the loaded data
contains three filers that genuinely list the subsidiary. All three arms were
naming a real parent and being scored wrong. Because every model answer is
stored, correcting the ground truth and re-scoring cost no additional inference;
it moved text-to-SQL from 28 to 29 passing attempts, vector RAG from 6 to 9 and
GraphRAG from 44 to 45. That we found it after believing the results were final
is direct evidence that others may remain.

\textbf{An asymmetry that penalised the graph, and what fixing it changed.}
The more instructive defect is one we initially reported as a finding. Both
stores keep a normalised name column, lowercased with corporate suffixes
stripped, so \emph{AEP Texas Inc.}\ is stored as \texttt{aep texas}. Neither
schema description said so. Text-to-SQL was unaffected because it matched the
raw column with \texttt{ILIKE}; the graph arm wrote an exact property match
against \texttt{nameNormalized} using the name as it appeared in the question,
retrieved zero rows, and refused. We first published that as a category loss for
structural retrieval.

It is not one. It is the same class of defect as an asymmetry we had already
found and fixed in the opposite direction --- an earlier iteration documented a
name convention to the graph but not to SQL, making text-to-SQL fail for a
reason unrelated to its architecture. Having recorded that, we should have
checked for its mirror image and did not.

We therefore documented the normalisation identically in both schema
descriptions, changing no retrieval code and no question, and re-ran all 180
attempts. \Cref{tab:beforeafter} reports both measurements. We present the
repaired numbers as the paper's result and the original alongside them, because
a fix that moves a headline in the author's favour should be visible rather than
quietly absorbed.

The general lesson is not that our harness had bugs. It is that a benchmark
harness is itself an instrument requiring calibration, and that a list of
defects already found is only useful if it is read as a checklist for the ones
not yet found.

\subsection{Artifacts}

The ingestion pipeline, the three retrieval arms, the scoring harness, the
question set with its ground truth, and the stored answers for all 180 attempts
are retained, which is what made the re-scoring in item 8 possible without new
inference. Every table and figure in this paper is generated programmatically from
the stored results rather than transcribed, so the numbers in the prose and the
numbers in the tables cannot drift apart.


\begin{widetab}[htbp]
\centering
\small
\begin{tabular}{@{}l cc cc cc@{}}
\toprule
& \multicolumn{2}{c}{\textbf{Text-to-SQL}} & \multicolumn{2}{c}{\textbf{Vector RAG}} & \multicolumn{2}{c}{\textbf{GraphRAG}} \\
\cmidrule(lr){2-3}\cmidrule(lr){4-5}\cmidrule(lr){6-7}
& pub. & rep. & pub. & rep. & pub. & rep. \\
\midrule
Questions solved (of 20) & 9 & \textbf{10}~\textcolor{okcol}{$\uparrow$} & 3 & \textcolor{gray!60}{--} & 14 & \textbf{16}~\textcolor{okcol}{$\uparrow$} \\
Attempts passed (of 60) & 29 & \textbf{36}~\textcolor{okcol}{$\uparrow$} & 9 & \textcolor{gray!60}{--} & 45 & \textbf{50}~\textcolor{okcol}{$\uparrow$} \\
Confident falsehoods & 2 & \textbf{6}~\textcolor{badcol}{$\uparrow$} & 0 & \textcolor{gray!60}{--} & 0 & \textcolor{gray!60}{--} \\
Correct but uncited & 16 & \textbf{10}~\textcolor{okcol}{$\downarrow$} & 6 & \textbf{12}~\textcolor{badcol}{$\uparrow$} & 0 & \textcolor{gray!60}{--} \\
Total tokens (thousands) & 158 & \textbf{144}~\textcolor{okcol}{$\downarrow$} & 119 & \textbf{118}~\textcolor{okcol}{$\downarrow$} & 780 & \textbf{521}~\textcolor{okcol}{$\downarrow$} \\
\midrule
T1 attempts (of 15) & 0 & \textbf{7}~\textcolor{okcol}{$\uparrow$} & 3 & \textcolor{gray!60}{--} & 15 & \textcolor{gray!60}{--} \\
T2 attempts (of 15) & 15 & \textcolor{gray!60}{--} & 0 & \textcolor{gray!60}{--} & 6 & \textcolor{gray!60}{--} \\
T3 attempts (of 15) & 0 & \textcolor{gray!60}{--} & 0 & \textcolor{gray!60}{--} & 15 & \textbf{14}~\textcolor{badcol}{$\downarrow$} \\
T4 attempts (of 15) & 14 & \textcolor{gray!60}{--} & 6 & \textcolor{gray!60}{--} & 9 & \textbf{15}~\textcolor{okcol}{$\uparrow$} \\
\bottomrule
\end{tabular}

\caption{The published measurement (\emph{pub.}) and the repaired one
(\emph{rep.}) after documenting the name normalisation identically in both
schema descriptions. Nothing else changed: same questions, same retrieval code,
same model and settings, and the statement timeout added later for the live
service was disabled so that exactly one variable differs. Shown in full because
the repair moves a headline in the author's favour.}
\label{tab:beforeafter}
\end{widetab}

\begin{widetab}[htbp]
\centering
\begin{tabular}{@{}llr@{}}
\toprule
\textbf{Resource} & \textbf{Specification} & \textbf{USD/mo} \\
\midrule
EC2 (app, graph, gateway) & m7g.xlarge, 4 vCPU / 16\,GB & 85.12 \\
EC2 storage               & 100\,GB gp3                 & 9.12 \\
RDS PostgreSQL            & db.t4g.medium               & 61.32 \\
RDS storage               & 100\,GB gp3                 & 9.12 \\
S3 + CloudFront           & demo volume                 & 0.50 \\
NAT gateway               & not used (S3 VPC endpoint)  & 0.00 \\
\midrule
\textbf{Total}            &                             & \textbf{165} \\
\bottomrule
\end{tabular}

\caption{Monthly infrastructure cost, on-demand pricing in
\texttt{ap-south-1}. The embedding pass used a temporarily resized 16-vCPU host
at 97 chunks/s; the instance was scaled back afterwards.}
\label{tab:cost}
\end{widetab}

% ════════════════════════════════════════════════════════════════════════
\section{Conclusion}
\label{sec:conclusion}

We measured three retrieval architectures over one corpus and one question set:
167{,}050 SEC filings from 10{,}511 companies, 508{,}714 embedded chunks and a
graph of 430{,}257 nodes and 561{,}144 edges, queried with 20 questions in four
shape categories, each asked three times, for 180 attempts scored
programmatically rather than by an LLM. GraphRAG solved 14 of 20 questions on all
three attempts, text-to-SQL 9 and vector RAG 3.

The aggregate is the least useful number in the paper. \textbf{GraphRAG lost two
of the four categories}: structured aggregation to text-to-SQL as predicted, at 24
times the token cost, and cross-document entity resolution to text-to-SQL against
prediction. The defensible summary is therefore not that graphs beat vectors but a
rule about question shape:

\begin{itemize}
  \item If the answer lies in one document and the entity is named, the binding
        constraint is \emph{entity filtering}, not ranking. Similarity search
        alone fails at scale --- 3 of 15 here --- and needs either a metadata
        filter on a stable identifier or a traversal to supply one.
  \item If the question is a \texttt{GROUP BY} over typed columns, use SQL. A
        traversal adds no information and adds a place to go wrong.
  \item If the question requires resolving \emph{who} an entity is, or joining
        facts that separate documents assert, structural retrieval is worth its
        cost. This is where the margin was unambiguous: 15 of 15 against 0 of 15
        for both baselines.
  \item If provenance must be auditable, the graph's advantage is categorical
        rather than marginal. It was never uncited across 60 attempts, because the
        traversal returns accession numbers before any text is read, whereas
        text-to-SQL was right-but-unciteable 16 times.
\end{itemize}

The trustworthiness finding is narrower but sharper than the accuracy finding.
Only 2 of 180 attempts asserted something verifiably false, and both were the same
error: text-to-SQL conflating distinct people who share a surname, matching on the
name the question supplied because that is what the text contains. Resolving
identity to an identifier before reading any text prevents that class of error
structurally rather than probabilistically. For an application where a confident
wrong answer costs more than a refusal, that property may matter more than the
pass rate.

We claim none of this beyond one corpus, one model and twenty questions.

\subsection*{Future work}

Four steps would materially strengthen these results. \emph{Re-run all 180
attempts per model}, turning the cross-model probe of \Cref{tab:crossmodel} into a
measurement and separating architectural effects from model effects. \emph{Repair
the entity-resolution layer} --- 0.12\% \texttt{:RESOLVES\_TO} coverage and the
name-normalization mismatch --- and re-measure T4, which is the only way to settle
whether that loss is architectural or incidental. \emph{Obtain an independent
question set}, since the present one was written by the author of the graph, which
is the study's most serious threat. And \emph{vary the retriever}: a longer-context
encoder with a metadata filter on CIK would test directly whether the T1 collapse
is a property of embeddings or of this particular configuration --- our own
diagnosis predicts the metadata filter, not the encoder, does the work.

